\chapter[KESIMPULAN DAN SARAN]{\\ KESIMPULAN DAN SARAN}

\section{Kesimpulan}
\begin{sectioncontent}
    \begin{enumerate}[leftmargin=0.75cm,label=\arabic*.]
        \item YOLOv8 Nano dan YOLO11 Nano berhasil diterapkan untuk klasifikasi penyakit daun tomat pada enam kelas, yaitu Tomato Early Blight, Tomato Healthy, Tomato Leaf Late Blight, Tomato Leaf Yellow Curl Virus, Tomato Mold Leaf, dan Tomato Septoria Leaf Spot. Dataset yang digunakan berjumlah 7.200 citra, dengan data uji sebanyak 720 citra yang terbagi seimbang pada setiap kelas.

        \item Berdasarkan hasil pengujian klasifikasi, YOLOv8 Nano memperoleh akurasi sebesar 88,47\%, sedangkan YOLO11 Nano memperoleh akurasi sebesar 88,06\%. Dari sisi metrik makro, YOLOv8 Nano memperoleh \textit{macro F1-score} sebesar 88,63\%, sedangkan YOLO11 Nano memperoleh 88,12\%. Hasil ini menunjukkan bahwa performa kedua model cukup berdekatan, tetapi YOLOv8 Nano sedikit lebih unggul pada skenario penelitian ini.

        \item Berdasarkan hasil pengujian pada Raspberry Pi 5, YOLOv8 Nano memperoleh rata-rata waktu inferensi sebesar 7,78 ms dengan FPS sebesar 128,52. Sementara itu, YOLO11 Nano memperoleh rata-rata waktu inferensi sebesar 7,86 ms dengan FPS sebesar 127,18. Dari sisi ukuran model dan penggunaan memori, YOLOv8 Nano juga sedikit lebih ringan dibandingkan YOLO11 Nano.

        \item Berdasarkan kombinasi hasil akurasi, \textit{macro F1-score}, waktu inferensi, FPS, ukuran model, dan penggunaan memori, YOLOv8 Nano menjadi model yang lebih direkomendasikan pada penelitian ini. Namun, rekomendasi tersebut berlaku pada dataset, konfigurasi pelatihan, format ONNX, dan perangkat Raspberry Pi 5 yang digunakan, sehingga hasilnya tidak dapat digeneralisasi secara mutlak untuk seluruh kondisi.
    \end{enumerate}
\end{sectioncontent}

\section{Saran}
\begin{sectioncontent}
    \hspace{\parindent}Berdasarkan hasil penelitian dan keterbatasan yang ditemukan, terdapat beberapa saran untuk pengembangan penelitian selanjutnya, yaitu sebagai berikut:

    \begin{enumerate}[leftmargin=0.75cm,label=\arabic*.]
        \item Penelitian selanjutnya dapat menggunakan dataset yang lebih bervariasi, terutama citra daun tomat yang diambil langsung dari kondisi lapangan dengan variasi pencahayaan, latar belakang, sudut pengambilan, jarak kamera, dan kualitas citra yang berbeda. Hal ini diperlukan untuk menguji kemampuan generalisasi model pada kondisi yang lebih nyata.

        \item Pengujian dapat dilakukan pada beberapa jenis perangkat \textit{edge} lain, seperti Raspberry Pi versi berbeda, NVIDIA Jetson Nano, Orange Pi, atau perangkat sejenis. Dengan begitu, performa model dapat dibandingkan pada perangkat dengan karakteristik komputasi yang berbeda.

        \item Penelitian selanjutnya dapat membandingkan beberapa format model hasil ekspor, seperti ONNX, NCNN, TensorFlow Lite, atau format lain yang mendukung perangkat \textit{edge}. Perbandingan ini penting karena runtime dan format model dapat memengaruhi waktu inferensi, penggunaan memori, dan stabilitas performa.

        \item Pengembangan penelitian dapat dilanjutkan ke tugas \textit{object detection} atau segmentasi citra agar sistem tidak hanya mampu menentukan jenis penyakit, tetapi juga dapat menunjukkan lokasi bagian daun yang terindikasi penyakit. Selain itu, sistem juga dapat dikembangkan menjadi aplikasi berbasis kamera agar lebih mendekati penggunaan langsung di lapangan.
    \end{enumerate}
\end{sectioncontent}